\subsection{Stars and Bars con Factoriales Precalculados (días de espera entre oasis)}
\label{alg:8-5}

\graybold{Preguntas guía:}

\begin{itemize}
\item ¿Debo distribuir una cantidad fija de elementos idénticos entre varios lugares?
\item ¿Cada lugar puede recibir cero o más elementos?
\item ¿El orden dentro de cada grupo no importa, únicamente cuántos elementos recibe cada uno?
\end{itemize}

\graybold{Utilidad:} Contar el número de formas de distribuir una cantidad fija de elementos indistinguibles entre varios grupos independientes. Es una de las aplicaciones más comunes de la técnica de \textit{Stars and Bars} (estrellas y barras), utilizada frecuentemente en problemas de combinatoria, programación competitiva y teoría de conteo.

\graybold{Complejidad:} Preprocesamiento de factoriales \bigO{\mathit{MAX}}; cada consulta \bigO{1}; espacio \bigO{\mathit{MAX}}.

\graybold{Fundamento teórico:} Si el viaje mínimo hasta el oasis destino requiere $O$ días y debe durar exactamente $T$ días, existen $extra=T-O$ días de espera que pueden distribuirse libremente antes de salir o después de cualquiera de los $D$ oasis recorridos, es decir, entre $D+1$ posiciones distintas. Como los días de espera son indistinguibles y cada posición puede recibir cualquier cantidad de ellos, el problema se reduce directamente al principio de \textit{Stars and Bars}, obteniendo un total de $\binom{extra+D}{D}$ planes posibles.

\graybold{Ejercicio aplicado}

Se conocen las posiciones de todos los oasis de una ruta. Para cada consulta se recibe un oasis destino y el número exacto de días que debe durar el viaje. Calcular cuántos planes distintos existen para distribuir los días de espera manteniendo siempre el mismo recorrido.

\graybold{I/O:} $5 \le N \le 20$, $1 \le M \le 10$. Las distancias a los oasis no superan los 60 días. Los factoriales se precomputan una única vez hasta MAX = 80 y cada consulta se responde en tiempo constante; ese tope alcanza mientras $T - O + D \le 80$ (por ejemplo, con $T \le 60$ y $D \le 20$). El código supone además $T \ge O$, es decir, que el tiempo pedido alcanza para llegar al destino.

\graybold{Código solución}

\begin{lstlisting}[language=Python]
import sys
input = sys.stdin.buffer.readline
cant, queries = map(int, input().split()) # numero de oasis y consultas
oasis = list(map(int, input().split())) # distancia minima hasta cada oasis

MAX = 80
factorial = [1] * (MAX + 1) # factoriales para responder combinaciones en O(1)

for i in range(1, MAX + 1):
    factorial[i] = factorial[i - 1] * i

def combination(n, k):
    return factorial[n] // (factorial[k] * factorial[n - k])

for _ in range(queries):
    destiny, time = map(int, input().split())
    extraTime = time - oasis[destiny - 1] # dias extra que deben distribuirse como esperas
    res = combination(extraTime + destiny, destiny) # Stars and Bars
    sys.stdout.write(str(res) + "\n")
\end{lstlisting}
